\subsection{李群的半直积}

设 I 为有限集，\((\mathbf{L}_{i})_{i\in I}\) 为一族李群。\(L=\prod_{i\in I}L_{i}\) 上的群结构和流形结构相容，因而 L 带有李群结构。L 称为这族李群 \((\mathbf{L}_{i})_{i\in I}\) 的积李群。

设 L、M 为李群，\(\sigma\) 为从 L 到群 M 的自同构群的同态。设 S 为关于 \(\sigma\) 的 L 与 M 的外半直积（Algebra, Chapter I, §6, no. 1, Definition 2）。

命题 7. 如果映射 \((m, l) \mapsto \sigma(l)m\) 从 \(\mathbf{M} \times \mathbf{L}\) 到 \(\mathbf{M}\) 是解析的，则群 \(\mathbf{S}\) 配以 \(\mathbf{M}\) 和 \(\mathbf{L}\) 的积流形结构后是李群。

对于 L 中的 \(l, l'\) 以及 M 中的 \(m, m'\)，

\[
\begin{array}{r l} (m, l) (m ^ {\prime}, l ^ {\prime}) ^ {- 1} & = m l l ^ {- 1} m ^ {\prime - 1} = m (\sigma (l l ^ {\prime - 1}) m ^ {\prime - 1}) l l ^ {\prime - 1} \\ & = (m (\sigma (l l ^ {\prime - 1}) m ^ {\prime - 1}), l l ^ {\prime - 1}) \end{array}
\]

从而得到命题。

如果命题 7 的条件成立，则称李群 S 为关于 \(\sigma\) 的 L 与 M 的（外）半直积李群。

显然，L（相应地，M）到 S 的典范嵌入是从 L（相应地，M）到 S 的一个李子群的同构，并将该李子群与 L（相应地，M）认同。S 到 L 的典范映射是李群态射。

反之，设 G 为李群，L、M 为两个李子群，并且群 G（代数上）是 L 与 M 的半直积（Algebra, Chapter I, § 6, no. 1）。记 \(\sigma(l)m = lml^{-1}\)，其中 \(l \in L\) 且 \(m \in M\)。则 \(\sigma\) 满足命题 7 的条件。因此我们可以构造关于 \(\sigma\) 的 L 与 M 的半直积李群 S。映射 \(j: (m, l) \mapsto ml\) 将 S 映到 G 上，是群同构且解析。若 \(j\) 是李群同构，则称李群 G 为 L 与 M 的（内）半直积，并认同 S 与 G。对于所有 \(g \in G\)，记 \(g = p(g)q(g)\)，其中 \(p(g) \in M\) 且 \(q(g) \in L\)。要使李群 G 成为 L 与 M 的半直积，充要条件是映射 \(p: G \to M\) 和 \(q: G \to L\) 中至少一个是解析的；此时二者都是解析的；或者，充要条件是 \(\mathrm{T}_{e}(\mathrm{G})\) 是 \(\mathrm{T}_{e}(\mathrm{M})\) 与 \(\mathrm{T}_{e}(\mathrm{L})\) 的拓扑直和（因为若此条件成立，\(j\) 在 \(e_{\mathrm{s}}\) 处是 étale 的）。

例。设 E 为赋范空间，\(G = \mathbf{GL}(E)\)，T 为 E 的平移群，A 为由 G 和 T 生成的 E 的置换群。群 A 在代数上是 G 与 T 的半直积。（如果 E 是有限维的，则 A 是 E 的仿射群，参见 Algebra, Chapter II, § 9, no. 4）。令 \(\sigma\) 为 G 在 E 上的恒等线性表示，S 为关于 \(\sigma\) 的 G 与 E 的外半直积。对于所有 \(x \in E\)，令 \(t_{x}\) 为由 x 定义的 E 的平移。映射 \((x, u) \mapsto t_{x} \circ u\) 是群 S 到群 A 的同构 \(\Phi\)。映射 \((x, u) \mapsto \sigma(u)x = u(x)\) 从 \(E \times \mathcal{L}(E)\) 到 E 是连续双线性的，因而是解析的；所以它限制到 \(E \times G\) 后也是解析的。因此，S 配以 E 和 G 的积流形结构后是李群。我们借助 \(\Phi\) 将此结构传到 A 上。于是 A 成为李群，即作为李群的 G 与 T 的内半直积。

命题 8. 设 G、H 为李群，\(p: \mathrm{G} \to \mathrm{H}\) 和 \(s: \mathrm{H} \to \mathrm{G}\) 为李群态射，满足 \(p \circ s = \mathrm{id}_{\mathrm{H}}\) 且 \(N = \operatorname{Ker} p\)。则 \(N\) 是 \(G\) 的李子群，\(s\) 是 \(H\) 到 \(G\) 的一个李子群的同构，并且李群 \(G\) 是 \(s(H)\) 与 \(N\) 的内半直积。

\(T_{e}(p) \circ T_{e}(s) = \mathrm{id}_{T_{e}(\mathrm{H})}\)，因此 \(p\)（相应地，\(s\)）是浸没（相应地，浸入）。由 Differentiable and Analytic Manifolds, R, 5.10.5，N 是 G 的李子群。另一方面，\(s\) 是 H 到 \(s(H)\) 的同胚，因此 \(s\) 是 H 到 G 的一个李子群的同构（Differentiable and Analytic Manifolds, R, 5.8.3）。最后，对于所有 \(g \in G\)，\(g = (s \circ p)(g)\)。且 \(n\) 是某个 \(n \in N\)；由于 \(s \circ p\) 是解析的，李群 G 是 \(s(H)\) 与 N 的半直积。

